\documentclass[11pt]{article}
\usepackage{icmlko}
\begin{document}

\title{쉬운 판만 샀다 --- 판을 제일 많이 올린 채택이 실무의 판에서는 음수다}
\author{월드모델 실험실}
\date{노트 536 · 2026년 8월 4일}
\maketitle

\begin{abstract}
노트 $528 \cdot 533$ 이 \emph{실무의 판은 두 배 차이}(후보 다섯 중 $1$--$2$ 등이
$1.97$ 배)임을 보였다. 그러면 물음은 하나다 --- \textbf{이 사이클의 채택 넷은
어느 배수 구간을 샀나.}

\begin{center}
\begin{tabular}{lrrrr}
\hline
채택 & \textbf{$<2$ 배} & 부호 & $2\!\sim\!5$ 배 & $5$ 배 이상 \\
\hline
\texttt{DOMDROP} & $\mathbf{-0.00016}$ & $3/8$ & $-0.00003$ & $\mathbf{+0.00316}$ \\
\texttt{TRAINCAP} & $+0.00066$ & $6/8$ & $+0.00071$ & $+0.00071$ \\
\texttt{TRAINW} & $+0.00052$ & $7/8$ & $+0.00020$ & $+0.00030$ \\
\texttt{TRAINR} & $\mathbf{+0.00122}$ & $7/8$ & $+0.00046$ & $+0.00258$ \\
\hline
\end{tabular}
\end{center}

사전 등록 ①(\texttt{TRAINR} 이 $<2$ 배에서 $1$ 등) \textbf{통과},
②($<2$ 배에서 $+0.005$ 이상 버는 채택이 있다) \textbf{반증} --- 최대가 기준의
$1/4$ 이다. \textbf{네 채택을 다 합쳐도 실무의 판은 $54.8\% \to 55.0\%$ 다.}
\end{abstract}

\begin{figure}[t]
\centering
\includegraphics[width=\textwidth]{figs/f1.pdf}
\caption{\textbf{가} --- 채택별 구간 이득. \textbf{나} --- 기준 정확도와 덮음.
\textbf{다} --- 읽기.}
\end{figure}

\section{판이 오른 것과 실무가 좋아진 것은 다른 이야기다}

이 사이클은 판을 $0.4812 \to 0.4932$ 로 올렸고 그것은 진짜다(노트 $526 \cdot
535$ 가 정의와 폭 양쪽에서 흔들어도 안 넘어졌다). 그런데 \emph{어디서} 벌었는지
쪼개 보면

\begin{center}
\begin{tabular}{lrr}
\hline
구간 & 기준 정확도 & 쌍 비중 \\
\hline
\textbf{$<2$ 배} (실무의 판) & $54.8\%$ & $32\%$ \\
$2\!\sim\!5$ 배 & $64.5\%$ & $29\%$ \\
$5$ 배 이상 & $80.6\%$ & $39\%$ \\
\hline
\end{tabular}
\end{center}

\begin{quote}
\textbf{판에서 제일 큰 채택인 \texttt{DOMDROP}($+0.0043$)이 $<2$ 배 구간에서는
음수다}($-0.00016$, $3/8$ 이라 $0$ 을 물지만 양수가 아니다). 그 이득의 거의
전부가 $5$ 배 이상 구간에서 나온다($+0.00316$).
\end{quote}

\section{\texttt{TRAINR} 만 다르다 --- 그러나 그것도 작다}

\texttt{TRAINR}(학습 쪽 최근성)은 $<2$ 배에서 $+0.00122$($7/8$)로 넷 중 유일하게
뚜렷하다. 노트 535 에서 \emph{표준화하면 이득이 $1.14$ 배로 커지는} 유일한
채택이었던 이유가 여기 있다 --- 표준화는 작은 배수 구간에 무게를 더 주고,
\texttt{TRAINR} 은 거기서 번다.

\begin{quote}
\textbf{그래도 $+0.00122$ 다.} 실무의 판 정확도 $54.8\%$ 를 $54.9\%$ 로 만드는
크기다. \textbf{접전을 사는 손잡이는 아직 없다.}
\end{quote}

\section{왜 하필 큰 배수인가 --- 가설}

\texttt{DOMDROP} 은 애니의 \emph{검색 추이} 축을 가린다. 검색 추이는 ``크게
흥한 것''을 잡는 신호일 것이다 --- 극단끼리의 비교($5$ 배 이상)에서는 강하게
작동하고 엎치락뒤치락하는 쌍에서는 아무 말도 안 한다.

이것은 \textbf{축 수준에서 검증할 수 있다} --- 계열을 하나씩 끄고 구간별 손해를
재면 ``어느 자료가 어느 판을 사는가''의 지도가 나온다. \textbf{노트 537 에서
잰다}(사전 등록 완료).

\section*{조항}

\begin{quote}
\textbf{채택을 적을 때 ``어느 배수 구간을 샀는지''를 함께 적는다.} 판이 오른
것만으로는 제품이 좋아졌다고 말할 수 없다 --- 이 사이클의 $+0.0120$ 중 실무의
판에 닿은 것은 $+0.002$ 도 안 된다.
\end{quote}

\section*{같은 물음이 세 번째로 돌아왔다}

노트 528 --- 관문의 자 다섯이 전부 전체 쌍이라 접전 개선을 못 본다.
노트 532 --- 신뢰도 신호가 접전에서 침묵한다.
여기 --- 채택 넷이 접전을 거의 안 산다.

\begin{quote}
\textbf{관문에 \emph{작은 배수 구간 전용 자}가 필요하다.} 지금 구조로는 접전을
사는 손잡이가 나와도 판이 그것을 거의 못 보고, 반대로 쉬운 판만 사는 손잡이가
판을 크게 올린다.
\end{quote}

\section*{남은 것}

\begin{itemize}
\item \textbf{노트 537} --- 어느 축 계열이 어느 구간을 사는가(진행 중).
\item \textbf{접전 전용 자를 관문에 넣을지} --- 넣으면 채택 넷을 그 자로 다시
심사해야 한다. \texttt{DOMDROP} 이 걸릴 수 있다.
\item \textbf{$5$ 배 이상 구간이 쌍의 $39\%$ 인데 실무 결정에는 거의 안
나온다} --- 판의 무게 자체가 실무와 어긋나 있다는 뜻이다.
\end{itemize}

\end{document}
